%% kapa: exam=1 round=32 year=2021 period=1 form=A subject=civil_law subject_ko=민법 seq=1 number=1
%% source: Q-Net 감정평가사 기출문제 (cst00309) artlSeq=5208397 file=extracted/1008/5208397/A형/2021년도 제32회 감정평가사 1차 1교시 A형.pdf
%% conversion: kordoc 4.24.51 -> tools/gen1.py (auto); review=proofread-sonnet; flags=none
\begin{question}{1}
관습법과 사실인 관습에 관한 설명으로 옳지 않은 것은? (다툼이 있으면 판례에 따름)
\begin{choices}
  \item 관습법은 성문법에 대하여 보충적 효력을 갖는다.
  \item 공동선조와 성과 본을 같이하는 미성년자인 후손은 종중의 구성원이 될 수 없다.
  \item 관습법이 성립한 후 사회구성원들이 그러한 관행의 법적 구속력에 더 이상 법적 확신을 갖지 않게 된 경우, 그 관습법은 법적 규범으로서의 효력이 없다.
  \item 사실인 관습은 법령으로서의 효력이 없고, 법률행위 당사자의 의사를 보충함에 그친다.
  \item 미등기 무허가건물의 매수인은 그 소유권이전등기를 마치지 않아도 소유권에 준하는 관습상의 물권을 취득한다.
\end{choices}
\end{question}
\answer{5}
